% Part: sets-functions-relations
% Chapter: infinite
% Section: dedekind-induction

\documentclass[../../../include/open-logic-section]{subfiles}

\begin{document}

\olfileid[id]{sfr}{infinite}{induction}
\olsection[Induksi Aritmetis]{Aljabar Dedekind dan Induksi Aritmetis}

Hal yang sangat penting ialah bahwa suatu aljabar Dedekind---bahkan,
\emph{setiap} aljabar Dedekind---dapat berfungsi sebagai pengganti bilangan
asli. Hal ini diperoleh berkat konsekuensi sederhana berikut:

\begin{thm}[Induksi aritmetis]\ollabel{thm:dedinfiniteinduction}
Misalkan $N, s, o$ membentuk suatu aljabar Dedekind. Maka, untuk sembarang
himpunan $X$:
	\begin{center}
		jika $o \in X$ dan $(\forall {n} \in N \cap X){s}({n}) \in X$, {maka} $N \subseteq X$.
	\end{center}
\end{thm}

\begin{proof}
Berdasarkan definisi aljabar Dedekind, $N = \closureofunder{s}{o}$.
Sekarang, jika ${o} \in X$ dan $(\forall {n} \in N)(n \in X \lif
{s}({n}) \in X)$, maka $N = \closureofunder{s}{o} \subseteq X$.
\end{proof}

Karena induksi merupakan ciri bilangan asli, pokoknya adalah sebagai berikut.
Dari sembarang himpunan yang tak hingga Dedekind, kita dapat membentuk suatu
aljabar Dedekind dan menggunakan aljabar tersebut sebagai pengganti bilangan
asli.

Memang, \olref{thm:dedinfiniteinduction} merumuskan induksi dalam istilah
\emph{teori himpunan}. Namun, prinsip itu mudah kita nyatakan dalam istilah
yang mungkin lebih akrab:

\begin{cor}\ollabel{natinductionschema}
Misalkan $N, s, o$ membentuk suatu aljabar Dedekind. Maka, untuk sembarang
rumus $\phi(x)$, yang boleh mempunyai parameter:
\begin{center}
	jika $\phi(o)$ dan $(\forall {n} \in N)(\phi(n)\lif
	\phi({s}({n})))$, {maka} $(\forall n \in N)\phi(n)$.
\end{center}
\end{cor}

\begin{proof}
Misalkan $X = \Setabs{n \in N}{\phi(n)}$, lalu gunakan
\olref{thm:dedinfiniteinduction}.
\end{proof}

Dalam hasil ini, kita berbicara tentang suatu rumus yang ``mempunyai
parameter''. Secara kasar, maksudnya ialah bahwa untuk sembarang objek
$c_1, \ldots, c_k$, kita dapat bekerja dengan
$\phi(x, c_1, \ldots, c_k)$. Secara lebih saksama, kita dapat menyatakan
hasil tersebut tanpa menyebut ``parameter'' sebagai berikut. Untuk sembarang
rumus $\phi(x, v_1, \ldots, v_k)$, yang semua variabel bebasnya ditampilkan,
kita mempunyai:
	\begin{align*}
			\forall v_1 \ldots \forall v_k((&\phi(o, v_1,\ldots, v_k) \land {}\\
			&	(\forall x \in N)(\phi(x,v_1, \ldots, v_k) \lif \phi(s(x), v_1,\ldots, v_k))) \lif {}\\
			&\hspace{3em} (\forall x \in N)\phi(x, v_1,\ldots, v_k))
	\end{align*}
Jelaslah bahwa pembicaraan mengenai ``mempunyai parameter'' dapat membuat
uraian jauh lebih mudah dibaca. (Dalam \olref[sth][][]{part}, kita akan cukup
sering menggunakan perangkat ini.)

Kembali ke aljabar Dedekind: dari sembarang aljabar Dedekind, kita juga dapat
mendefinisikan fungsi-fungsi aritmetis yang biasa, yakni penjumlahan,
perkalian, dan perpangkatan. Namun, hal ini tidak sepele dan melibatkan
teknik \emph{definisi rekursif}. Teknik tersebut akan kita perkenalkan dan
benarkan jauh kemudian, dan dalam konteks yang jauh lebih umum. (Pembaca
yang berminat mungkin ingin meninjau kembali hal ini setelah
\olref[sth][ord-arithmetic][]{chap}, atau mungkin membaca pembahasan
alternatif, misalnya \citealt[pp.~95--8]{Potter2004}.) Namun, jika
$N, s, o$ membentuk suatu aljabar Dedekind, pada akhirnya kita akan dapat
menetapkan hal-hal berikut:
\begin{align*}
	{a} + {o} &= {a} & & & {a} \times {o} &= {o} & & & {a}^{o} &= s(o)\\
{a} + {s}({b}) &= {s}({a}+{b}) &&& {a} \times {s}({b}) &= ({a}\times {b}) + {a}  & & & {a}^{{s}({b})} &= {a}^{b} \times {a}
\end{align*}
dan menunjukkan bahwa semuanya berperilaku sebagaimana yang diharapkan.

\end{document}
